\documentclass[12pt, titlepage]{article}

\usepackage{booktabs}
\usepackage{tabularx}
\usepackage{hyperref}
\hypersetup{
    colorlinks,
    citecolor=blue,
    filecolor=black,
    linkcolor=red,
    urlcolor=blue
}
\usepackage[round]{natbib}

\input{../Comments.text}
\input{../Common.text}

\begin{document}

\title{System Verification and Validation Plan for \progname{}} 
\author{\authname}
\date{\today}
	
\maketitle

\pagenumbering{roman}

\section*{Revision History}

\begin{tabularx}{\textwidth}{p{3cm}p{2cm}X}
\toprule {\bf Date} & {\bf Version} & {\bf Notes}\\
\midrule
Date 1 & 1.0 & Notes\\
Date 2 & 1.1 & Notes\\
\bottomrule
\end{tabularx}

~\\
\wss{The intention of the VnV plan is to increase confidence in the software.
However, this does not mean listing every verification and validation technique
that has ever been devised.  The VnV plan should also be a \textbf{feasible}
plan. Execution of the plan should be possible with the time and team available.
If the full plan cannot be completed during the time available, it can either be
modified to ``fake it'', or a better solution is to add a section describing
what work has been completed and what work is still planned for the future.}

\wss{The VnV plan is typically started after the requirements stage, but before
the design stage.  This means that the sections related to unit testing cannot
initially be completed.  The sections will be filled in after the design stage
is complete.  the final version of the VnV plan should have all sections filled
in.}

\newpage

\tableofcontents

\listoftables
\wss{Remove this section if it isn't needed}

\listoffigures
\wss{Remove this section if it isn't needed}

\newpage

\section{Symbols, Abbreviations, and Acronyms}

\renewcommand{\arraystretch}{1.2}
\begin{tabular}{l l} 
  \toprule		
  \textbf{symbol} & \textbf{description}\\
  \midrule 
  T & Test\\
  \bottomrule
\end{tabular}\\

\wss{symbols, abbreviations, or acronyms --- you can simply reference the SRS
  \citep{SRS} tables, if appropriate}

\wss{Remove this section if it isn't needed}

\newpage

\pagenumbering{arabic}

This document ... \wss{provide an introductory blurb and roadmap of the
  Verification and Validation plan}

\section{General Information}

\subsection{Summary}

\wss{Say what software is being tested.  Give its name and a brief overview of
  its general functions.}

\subsection{Objectives}

\wss{State what is intended to be accomplished.  The objective will be around
  the qualities that are most important for your project.  You might have
  something like: ``build confidence in the software correctness,''
  ``demonstrate adequate usability.'' etc.  You won't list all of the qualities,
  just those that are most important.}

\wss{You should also list the objectives that are out of scope.  You don't have 
the resources to do everything, so what will you be leaving out.  For instance, 
if you are not going to verify the quality of usability, state this.  It is also 
worthwhile to justify why the objectives are left out.}

\wss{The objectives are important because they highlight that you are aware of 
limitations in your resources for verification and validation.  You can't do everything, 
so what are you going to prioritize?  As an example, if your system depends on an 
external library, you can explicitly state that you will assume that external library 
has already been verified by its implementation team.}

\subsection{Custom Documents}

\wss{Summarize the CDs (Custom Documents) that were tackled by this project.  CDs 
can include usability testing, code walkthroughs, user documentation, formal
proof, GenderMag personas, Design Thinking, etc.  CDs should have already
been approved by the course instructor as included in your problem statement.
You can use a pull request to update your CDs (in TeamComposition.csv or
Repos.csv) if your plan changes as a result of the VnV planning exercise.}

\subsection{Relevant Documentation}

\wss{Reference relevant documentation.  This will definitely include your SRS
  and your other project documents (design documents, like MG, MIS, etc).  You
  can include these even before they are written, since by the time the project
  is done, they will be written.  You can create BibTeX entries for your
  documents and within those entries include a hyperlink to the documents.}

\citet{SRS}

\wss{Don't just list the other documents.  You should explain why they are relevant and 
how they relate to your VnV efforts.}

\section{Plan}

\wss{Introduce this section.  You can provide a roadmap of the sections to
  come.}

\subsection{Verification and Validation Team}

\wss{Your teammates.  Maybe your supervisor.
  You should do more than list names.  You should say what each person's role is
  for the project's verification.  A table is a good way to summarize this information.}

\subsection{SRS Verification}

\wss{List any approaches you intend to use for SRS verification.  This may
  include ad hoc feedback from reviewers, like your classmates (like your
  primary reviewer), or you may plan for something more rigorous/systematic.}

\wss{If you have a supervisor for the project, you shouldn't just say they will
read over the SRS.  You should explain your structured approach to the review.
Will you have a meeting?  What will you present?  What questions will you ask?
Will you give them instructions for a task-based inspection?  Will you use your
issue tracker?}

\wss{Maybe create an SRS checklist?}

\subsection{Design Verification}

\wss{Plans for design verification}

\wss{The review will include reviews by your classmates}

\wss{Create a checklists?}

\subsection{Verification and Validation Plan Verification}

\wss{The verification and validation plan is an artifact that should also be
verified.  Techniques for this include review and mutation testing.}

\wss{The review will include reviews by your classmates}

\wss{Create a checklists?}

\subsection{Implementation Verification}

\wss{You should at least point to the tests listed in this document and the unit
  testing plan.}

\wss{In this section you would also give any details of any plans for static
  verification of the implementation.  Potential techniques include code
  walkthroughs, code inspection, static analyzers, etc.}

\wss{The final class presentation in CAS 741 could be used as a code
walkthrough.  There is also a possibility of using the final presentation (in
CAS741) for a partial usability survey.}

\subsection{Automated Testing and Verification Tools}

\wss{What tools are you using for automated testing.  Likely a unit testing
  framework and maybe a profiling tool, like ValGrind.  Other possible tools
  include a static analyzer, make, continuous integration tools, test coverage
  tools, etc.  Explain your plans for summarizing code coverage metrics.
  Linters are another important class of tools.  For the programming language
  you select, you should look at the available linters.  There may also be tools
  that verify that coding standards have been respected, like flake9 for
  Python.}

\wss{If you have already done this in the development plan, you can point to
that document.}

\wss{The details of this section will likely evolve as you get closer to the
  implementation.}

\subsection{Software Validation}

\wss{If there is any external data that can be used for validation, you should
  point to it here.  If there are no plans for validation, you should state that
  here.}

\wss{You might want to use review sessions with the stakeholder to check that
the requirements document captures the right requirements.  Maybe task based
inspection?}

\wss{For those capstone teams with an external supervisor, the Rev 0 demo should 
be used as an opportunity to validate the requirements.  You should plan on 
demonstrating your project to your supervisor shortly after the scheduled Rev 0 demo.  
The feedback from your supervisor will be very useful for improving your project.}

\wss{For teams without an external supervisor, user testing can serve the same purpose 
as a Rev 0 demo for the supervisor.}

\wss{This section might reference back to the SRS verification section.}

\section{System Tests}

\wss{There should be text between all headings, even if it is just a roadmap of
the contents of the subsections.}

\subsection{Tests for Functional Requirements}

The subsections below group the functional requirements in the Functional Requirements section of the
SRS \citep{SRS} by feature area. Together they cover FUNC-01 to FUNC-16.
FUNC-17 (Sponsored Reward Redemption) is a stretch goal in the SRS Waiting
Room, so it is not tested.

\wss{Each test uses the fields Control, Initial State, Input, Output, Test
Case Derivation, and How test will be performed.}

\subsubsection{Account Management}

These tests cover creating an account and signing in.

\noindent\begin{tabularx}{\textwidth}{p{3cm}X}
  \toprule
  \textbf{SRS ID} & \textbf{Requirement} \\
  \midrule
  FUNC-01 & User Registration \\
  FUNC-02 & User Authentication \\
  \bottomrule
\end{tabularx}

\wss{Add the test cases for these requirements.}

\subsubsection{Fitness Tracking and Goals}

These tests cover importing activity data from third-party providers and creating and tracking fitness goals.

\noindent\begin{tabularx}{\textwidth}{p{3cm}X}
  \toprule
  \textbf{SRS ID} & \textbf{Requirement} \\
  \midrule
  FUNC-03 & Third-Party Fitness Data Import \\
  FUNC-04 & Fitness Goal Creation \\
  FUNC-05 & Fitness Goal Tracking \\
  \bottomrule
\end{tabularx}

\wss{Add the test cases for these requirements.}

\subsubsection{Rewards and Virtual Currency}

These tests cover awarding virtual currency for activity and keeping each user's balance.

\noindent\begin{tabularx}{\textwidth}{p{3cm}X}
  \toprule
  \textbf{SRS ID} & \textbf{Requirement} \\
  \midrule
  FUNC-06 & Reward Allocation \\
  FUNC-07 & Virtual Currency Tracking \\
  \bottomrule
\end{tabularx}

\wss{Add the test cases for these requirements.}

\subsubsection{Simulated Investing}

These tests cover the portfolio and trading features of the simulated market.

\noindent\begin{tabularx}{\textwidth}{p{3cm}X}
  \toprule
  \textbf{SRS ID} & \textbf{Requirement} \\
  \midrule
  FUNC-08 & Portfolio Creation \\
  FUNC-09 & Asset Purchase \\
  FUNC-10 & Asset Sale \\
  FUNC-11 & Transaction Recording \\
  FUNC-12 & Portfolio Performance Calculation \\
  \bottomrule
\end{tabularx}

\wss{Add the test cases for these requirements.}

\subsubsection{Community and Leaderboards}

These tests cover friends lists and public and private leaderboard groups.

\noindent\begin{tabularx}{\textwidth}{p{3cm}X}
  \toprule
  \textbf{SRS ID} & \textbf{Requirement} \\
  \midrule
  FUNC-13 & Leaderboard Participation \\
  FUNC-14 & Friend Management \\
  FUNC-15 & Private Leaderboards \\
  \bottomrule
\end{tabularx}

\wss{Add the test cases for these requirements.}

\subsubsection{Progress Reporting}

These tests cover the summaries of activity, goals, earnings, portfolio performance, and standings.

\noindent\begin{tabularx}{\textwidth}{p{3cm}X}
  \toprule
  \textbf{SRS ID} & \textbf{Requirement} \\
  \midrule
  FUNC-16 & Progress Reporting \\
  \bottomrule
\end{tabularx}

\wss{Add the test cases for these requirements.}

\subsection{Tests for Nonfunctional Requirements}

The subsections below follow the nonfunctional requirement sections of the
SRS \citep{SRS}. Each one lists the SRS subsections and requirement IDs that
its tests cover.

\wss{Each test uses the fields Type, Initial State, Input/Condition,
Output/Result, and How test will be performed. Static tests, such as reviews,
inspections, and walkthroughs, should instead say who is involved and how the
review is run.}

\subsubsection{Look and Feel}

These tests cover the Look and Feel Requirements section of the SRS.

\paragraph{Appearance}

Covers LF-APP-01 to LF-APP-05.

\wss{Add the test cases for these requirements.}

\paragraph{Style}

Covers LF-STY-01 to LF-STY-04.

\wss{Add the test cases for these requirements.}

\subsubsection{Usability and Humanity}

These tests cover the Usability and Humanity Requirements section of the SRS.

\paragraph{Ease of Use}

Covers UHR-EOU-01 to UHR-EOU-09.

\wss{Add the test cases for these requirements.}

\paragraph{Personalization and Internationalization}

Covers UHR-PER-01 to UHR-PER-04, UHR-INT-01.

\wss{Add the test cases for these requirements.}

\paragraph{Learning}

Covers UHR-LRN-01 to UHR-LRN-04.

\wss{Add the test cases for these requirements.}

\paragraph{Understandability and Politeness}

Covers UHR-UND-01 to UHR-UND-03, UHR-POL-01 to UHR-POL-05.

\wss{Add the test cases for these requirements.}

\paragraph{Accessibility}

Covers UHR-ACC-01 to UHR-ACC-06.

\wss{Add the test cases for these requirements.}

\subsubsection{Performance}

These tests cover the Performance Requirements section of the SRS.

\paragraph{Speed and Latency}

Covers PR-SL-01 to PR-SL-04.

\wss{Add the test cases for these requirements.}

\paragraph{Safety-Critical}

Covers PR-SC-01 to PR-SC-08.

\wss{Add the test cases for these requirements.}

\paragraph{Precision or Accuracy}

Covers PR-PA-01 to PR-PA-06.

\wss{Add the test cases for these requirements.}

\paragraph{Robustness or Fault-Tolerance}

Covers PR-RFT-01 to PR-RFT-04.

\wss{Add the test cases for these requirements.}

\paragraph{Capacity}

Covers PR-C-01, PR-C-02.

\wss{Add the test cases for these requirements.}

\paragraph{Scalability or Extensibility}

Covers PR-SE-01 to PR-SE-03.

\wss{Add the test cases for these requirements.}

\paragraph{Longevity}

Covers PR-L-01, PR-L-02.

\wss{Add the test cases for these requirements.}

\subsubsection{Operational and Environmental}

These tests cover the Operational and Environmental Requirements section of the SRS.

\paragraph{Productization}

Covers PROD-01, PROD-02.

\wss{Add the test cases for these requirements.}

\subsubsection{Maintainability and Support}

These tests cover the Maintainability and Support Requirements section of the SRS.

\paragraph{Maintenance}

Covers MAS-MNT-01 to MAS-MNT-05.

\wss{Add the test cases for these requirements.}

\paragraph{Supportability}

Covers MAS-SUP-01.

\wss{Add the test cases for these requirements.}

\paragraph{Adaptability}

Covers MAS-ADPT-01 to MAS-ADPT-03.

\wss{Add the test cases for these requirements.}

\subsubsection{Security}

These tests cover the Security Requirements section of the SRS \citep{SRS}, spanning access control, data and logic integrity, privacy, economic and security audit logging, and immunity against attacks.

\paragraph{Access}

Covers SEC-ACC-01 to SEC-ACC-05.

\begin{enumerate}

\item{\textbf{T-SEC-ACC-01: Authenticated Access Token Validation}\\}

Type: Dynamic, Automated

Initial State: Backend server running in test profile with test database; user account registered; valid access token issued with a 60-minute expiry window.

Input/Condition: Send requests to a protected system resource (such as fitness activities or portfolio summaries) under four conditions: (1) missing authentication credentials, (2) tampered token signature, (3) expired token (timestamp beyond the 60-minute lifetime), and (4) authentic, unexpired token.

Output/Result: Requests lacking credentials, with a tampered token, or with an expired token are rejected with HTTP 401 Unauthorized. The request with a valid token returns HTTP 200 OK with the requested user data.

How test will be performed: Automated integration tests dispatching requests across token states against protected system services.

\item{\textbf{T-SEC-ACC-02: Role-Based Authorization and Cross-User Isolation}\\}

Type: Dynamic, Automated

Initial State: Two standard test users (User A and User B) with distinct portfolios and data records, and one administrator user.

Input/Condition: (1) User A attempts to read User B's private portfolio. (2) User A attempts to execute a trade or modify records on User B's behalf. (3) User A attempts to access system-wide administrative audit logs. (4) Administrator requests access to system-wide administrative audit logs.

Output/Result: Cross-user reads, cross-user modifications, and standard-user administrative requests are rejected with HTTP 403 Forbidden with zero data modifications. The administrative request is accepted (HTTP 200 OK) returning the requested audit records.

How test will be performed: Automated security integration tests verifying role-based access control rules and tenant isolation across system boundaries.

\item{\textbf{T-SEC-ACC-03: Secure Credential Storage and Secret Exclusion}\\}

Type: Static and Dynamic, Automated

Initial State: Repository source commit under evaluation; mobile test environment.

Input/Condition: (1) Execute secret scanning on the complete version control repository history for API keys, private keys, and client secrets. (2) Store third-party fitness integration credentials via the client application and inspect local persistent device storage.

Output/Result: (1) Repository secret scan detects zero exposed credentials or keys. (2) Integration credentials are encrypted at rest using platform-managed secure key storage (e.g., AES-256 GCM); zero cleartext tokens exist on local disk.

How test will be performed: Automated repository secret scanning in CI and mobile client instrumentation tests inspecting device storage.

\item{\textbf{T-SEC-ACC-04: Password Strength, Lockout, and Multi-Factor Authentication}\\}

Type: Dynamic, Automated

Initial State: User authentication service active with test database; common password dictionary loaded in memory.

Input/Condition: (1) Submit account registration requests using short passwords ($< 8$ characters), common dictionary passwords, and a strong password ($> 8$ characters). (2) Submit 5 consecutive incorrect password attempts for a user account, followed immediately by a 6th attempt with the correct password within 15 minutes. (3) Advance system time past 15 minutes (or submit a password reset) and retry with valid credentials. (4) For an account with multi-factor authentication (MFA) enabled, submit valid credentials, followed by an invalid one-time password (OTP), and then a valid OTP.

Output/Result: Short and dictionary passwords return HTTP 400 Bad Request; strong passwords return HTTP 201 Created. The 6th consecutive failed attempt triggers account lockout, rejecting subsequent login attempts with HTTP 423 Locked (or HTTP 403 Forbidden). Login succeeds once the lockout duration elapses. MFA challenge rejects invalid OTPs (HTTP 401 Unauthorized) and grants session access only upon valid OTP submission.

How test will be performed: Automated integration tests verifying authentication policies, lockout timers, and MFA verification flows.

\item{\textbf{T-SEC-ACC-05: Session Termination and Device Unlinking}\\}

Type: Dynamic, Automated

Initial State: Active user session with valid access and refresh tokens; client device linked to user account.

Input/Condition: (1) Client issues a sign-out request with its registered device identifier. (2) Attempt subsequent resource requests using the revoked access token. (3) Attempt to submit activity data from the unlinked device. (4) Present a session token that has been inactive for greater than 30 days.

Output/Result: Sign-out completes successfully (HTTP 200 OK) and unlinks the device. Requests presenting revoked tokens, data submissions from unlinked devices, and expired inactive sessions are rejected (HTTP 401 Unauthorized or HTTP 403 Forbidden).

How test will be performed: Automated integration tests verifying session invalidation, token revocation repositories, and device link states.

\end{enumerate}

\paragraph{Integrity}

Covers SEC-INT-01 to SEC-INT-03.

\begin{enumerate}

\item{\textbf{T-SEC-INT-01: Activity Telemetry Plausibility and Deduplication}\\}

Type: Dynamic, Automated

Initial State: Authenticated user account with 0 virtual currency.

Input/Condition: Ingest four workout telemetry payloads: (1) sustained cadence of 280 steps/min (exceeding plausibility threshold of 250 steps/min), (2) overlapping workout intervals from two integrations, (3) duplicate submission of an identical workout external sync ID, and (4) normal plausible workout (120 steps/min, no overlap).

Output/Result: (1) Implausible cadence payload marked \texttt{UNVERIFIED} with zero virtual currency awarded. (2) Overlapping interval flagged \texttt{UNVERIFIED}; only unique segments credited. (3) Duplicate payload rejected by idempotency key with zero duplicate currency minted. (4) Plausible workout verified and virtual currency awarded.

How test will be performed: Automated unit and integration tests executing domain validation rules on activity telemetry payloads.

\item{\textbf{T-SEC-INT-02: Server-Side Authority on Valuations and Trades}\\}

Type: Dynamic, Automated

Initial State: Authenticated user with virtual currency balance; server-authoritative asset price set (e.g., 50.00 FLX for \texttt{FLEX\_CORP}).

Input/Condition: (1) Client submits a trade request specifying a client-dictated execution price (e.g., 0.01 FLX) and total cost. (2) Client submits a currency reward claim specifying an arbitrary client-dictated reward amount.

Output/Result: (1) The system rejects the trade request with HTTP 400 Bad Request, or ignores the client-supplied price and executes strictly at the server-authoritative valuation (50.00 FLX). (2) The system rejects client-dictated reward amounts and calculates currency rewards strictly from verified workout metrics.

How test will be performed: Automated integration tests verifying server-side pricing enforcement and reward calculation authority.

\item{\textbf{T-SEC-INT-03: Database Backup Integrity and Restoration Drill}\\}

Type: Dynamic Operational Verification and Static Configuration Review

Initial State: Staging database populated with seed dataset of users, portfolios, and audit records; automated backup cron job configured.

Input/Condition: (1) Inspect cloud backup policy scripts for 24-hour backup interval and 30-day retention lifecycle. (2) Trigger an automated backup snapshot, record SHA-256 data checksums, restore onto a clean containerized database instance, and measure restoration time.

Output/Result: Configuration confirms daily backups and 30-day automated retention. Database restores completely within target limit ($< 15$ minutes in test environment), with record counts and checksums matching 100\%.

How test will be performed: Automated staging script executing backup creation, containerized restoration, and data integrity verification.

\end{enumerate}

\paragraph{Privacy}

Covers SEC-PRV-01 to SEC-PRV-03.

\begin{enumerate}

\item{\textbf{T-SEC-PRV-01: Encryption in Transit and at Rest}\\}

Type: Dynamic Security Verification and Database Inspection

Initial State: Staging backend deployed behind reverse proxy; database running with encrypted tablespace.

Input/Condition: (1) Execute network handshake scan against API endpoints testing plain HTTP and TLS protocols 1.0, 1.1, 1.2, and 1.3. (2) Insert sensitive biometric fields (heart rate, body weight) into database and inspect raw storage files on disk.

Output/Result: (1) Plain HTTP requests and TLS connections below version 1.3 are rejected; only TLS 1.3 handshakes succeed. (2) Direct inspection of database disk blocks reveals AES-256 ciphertext; sensitive biometric data is unreadable in raw form.

How test will be performed: Automated script combining TLS protocol scanning and raw disk inspection on test database volumes.

\item{\textbf{T-SEC-PRV-02: Private by Default and Leaderboard Opt-In Isolation}\\}

Type: Dynamic, Automated

Initial State: Newly registered user with logged workouts and simulated portfolio; default settings unchanged.

Input/Condition: (1) Query public user profile and public leaderboard for the new user. (2) User explicitly opts in to leaderboard sharing with a public alias. (3) User toggles leaderboard sharing off.

Output/Result: (1) By default, public profile returns 404 or minimal stub; user is omitted from public leaderboards. (2) Opted-in leaderboard displays only chosen alias and aggregate score; real name, email, and biometric metrics remain private. (3) Opting out removes the user from public leaderboards immediately.

How test will be performed: Automated API integration tests verifying endpoint responses across default, opt-in, and opt-out states.

\item{\textbf{T-SEC-PRV-03: User Account Deletion and PII Eradication}\\}

Type: Dynamic, Automated

Initial State: User account populated with profile, health logs, active session, and transactions in the financial audit ledger.

Input/Condition: User submits an account deletion request with credential confirmation. Inspect active application databases, attempt login with deleted credentials, and inspect the financial audit ledger.

Output/Result: All personal identifying information (name, email, device identifiers) and health logs are purged from the active system within 24 hours. Subsequent login attempts are rejected (HTTP 401 Unauthorized). The financial audit ledger preserves transactional continuity but replaces user identifiers with an irreversible pseudonym.

How test will be performed: Automated end-to-end integration test asserting database purge, login denial, and ledger anonymization.

\end{enumerate}

\paragraph{Audit}

Covers SEC-AUD-01, SEC-AUD-02.

\begin{enumerate}

\item{\textbf{T-SEC-AUD-01: Economic Event Ledger Immutability}\\}

Type: Dynamic Database Privilege Verification

Initial State: Backend connected to database using standard application credentials; sample rows present in \texttt{economic\_audit\_ledger}.

Input/Condition: (1) Record simulated trade or reward via standard service API. (2) Attempt direct SQL \texttt{UPDATE} on an existing ledger record. (3) Attempt direct SQL \texttt{DELETE} on an existing ledger record.

Output/Result: Service API inserts new audit record with immutable timestamp. Direct SQL \texttt{UPDATE} and \texttt{DELETE} commands are denied by database role permissions with error 1142 (command denied); ledger records remain unmodified.

How test will be performed: Automated database integration test executing restricted SQL commands using application service credentials.

\item{\textbf{T-SEC-AUD-02: Structured Security Event Logging}\\}

Type: Dynamic, Automated

Initial State: Log appender active and configured for structured JSON output.

Input/Condition: Trigger security events: 3 failed login attempts, activity payload with tampered signature, rate limit exceeded ($> 60$ req/min), and unhandled server exception.

Output/Result: Logs capture structured JSON entries containing severity level (WARN/ERROR), ISO 8601 UTC timestamp, client IP, event type, and user ID (or anonymous). Zero sensitive credentials or plaintext passwords appear in log records.

How test will be performed: Automated integration test with in-memory log appender asserting on serialized JSON log structure and contents.

\end{enumerate}

\paragraph{Immunity}

Covers SEC-IMM-01 to SEC-IMM-03.

\begin{enumerate}

\item{\textbf{T-SEC-IMM-01: Injection Defense (SQLi, XSS, Command Injection)}\\}

Type: Static SAST and Dynamic DAST Penetration Testing

Initial State: Backend service deployed in containerized test environment with database connectivity.

Input/Condition: (1) Run static analysis on system source code for unparameterized queries, unescaped reflection, and command execution. (2) Execute automated penetration scanning dispatching SQL injection, cross-site scripting (XSS), and command injection attack vectors across all external system interfaces.

Output/Result: Static analysis reports zero high-severity injection vulnerabilities. Penetration scanning reports zero exploitable injection vulnerabilities; malicious payloads are sanitized or rejected (HTTP 400/422).

How test will be performed: Automated security test pipeline combining static code analysis (SAST) and dynamic penetration scanning (DAST).

\item{\textbf{T-SEC-IMM-02: API Rate Limiting Enforcement}\\}

Type: Dynamic, Automated

Initial State: Rate limiting filter configured with threshold of 60 requests/minute per client IP.

Input/Condition: Client sends 60 consecutive requests within 30 seconds, sends 10 additional requests within the remainder of the 60-second window, waits 65 seconds, and sends a subsequent request.

Output/Result: Requests 1 to 60 process normally. Requests 61 to 70 return HTTP 429 Too Many Requests with a \texttt{Retry-After} header. The request following the cooldown period processes normally.

How test will be performed: Automated integration test using a multi-threaded HTTP client dispatching timed request bursts.

\item{\textbf{T-SEC-IMM-03: Client Environment Tampering and Root Detection}\\}

Type: Dynamic Verification

Initial State: Android test device or emulator with root privileges enabled (\texttt{su} binary present, test keys installed).

Input/Condition: (1) Deploy production release build onto rooted device and attempt app launch, fitness sync, and leaderboard access. (2) Deploy debug build onto rooted device and attempt same actions.

Output/Result: Production build detects untrusted environment, displays warning modal, and disables virtual currency minting and leaderboard features. Debug build runs without restrictions for testing.

How test will be performed: Automated Android UI test and manual release checklist verifying root detection library signals.

\end{enumerate}

\paragraph{Security Review Walkthrough}

Static reviews are conducted once per milestone release cycle (Rev 0, Final Demo) by the Backend Lead (Daniel Franze-Da Silva), Security and Documentation Lead (Arjun Karthik), and an independent peer reviewer. The team reviews: (1) ISO/IEC 27001 Annex A mapping against system controls (COMP-STD-07), (2) least privilege scopes for third-party fitness APIs (SEC-ACC-03), (3) repository secret scan results, and (4) dependency vulnerability reports ensuring zero unresolved High/Critical CVEs.

\subsubsection{Cultural}

These tests cover the Cultural Requirements section of the SRS.

\paragraph{Cultural}

Covers CUL-01, CUL-02.

\wss{Add the test cases for these requirements.}

\subsubsection{Compliance}

These tests cover the Compliance Requirements section of the SRS.

\paragraph{Legal}

Covers COMP-LEG-01 to COMP-LEG-07.

\wss{Add the test cases for these requirements.}

\paragraph{Standards Compliance}

Covers COMP-STD-01 to COMP-STD-07.

\wss{Add the test cases for these requirements.}

\subsubsection{User Documentation and Training}

These tests cover the User Documentation and Training Requirements section of the SRS.

\paragraph{User Documentation}

Covers DOC-01 to DOC-04.

\wss{Add the test cases for these requirements.}

\paragraph{Training}

Covers TRN-01 to TRN-03.

\wss{Add the test cases for these requirements.}

\subsection{Traceability Between Test Cases and Requirements}

\wss{Provide a table that shows which test cases are supporting which
  requirements.}

\section{Unit Test Description}

\wss{This section should not be filled in until after the MIS (detailed design
  document) has been completed.}

\wss{Reference your MIS (detailed design document) and explain your overall
philosophy for test case selection.}  

\wss{To save space and time, it may be an option to provide less detail in this section.  
For the unit tests you can potentially layout your testing strategy here.  That is, you 
can explain how tests will be selected for each module.  For instance, your test building 
approach could be test cases for each access program, including one test for normal behaviour 
and as many tests as needed for edge cases.  Rather than create the details of the input 
and output here, you could point to the unit testing code.  For this to work, you code 
needs to be well-documented, with meaningful names for all of the tests.}

\subsection{Unit Testing Scope}

\wss{What modules are outside of the scope.  If there are modules that are
  developed by someone else, then you would say here if you aren't planning on
  verifying them.  There may also be modules that are part of your software, but
  have a lower priority for verification than others.  If this is the case,
  explain your rationale for the ranking of module importance.}

\subsection{Tests for Functional Requirements}

\wss{Most of the verification will be through automated unit testing.  If
  appropriate specific modules can be verified by a non-testing based
  technique.  That can also be documented in this section.}

\subsubsection{Module 1}

\wss{Include a blurb here to explain why the subsections below cover the module.
  References to the MIS would be good.  You will want tests from a black box
  perspective and from a white box perspective.  Explain to the reader how the
  tests were selected.}

\begin{enumerate}

\item{test-id1\\}

Type: \wss{Functional, Dynamic, Manual, Automatic, Static etc. Most will
  be automatic}
					
Initial State: 
					
Input: 
					
Output: \wss{The expected result for the given inputs}

Test Case Derivation: \wss{Justify the expected value given in the Output field}

How test will be performed: 
					
\item{test-id2\\}

Type: \wss{Functional, Dynamic, Manual, Automatic, Static etc. Most will
  be automatic}
					
Initial State: 
					
Input: 
					
Output: \wss{The expected result for the given inputs}

Test Case Derivation: \wss{Justify the expected value given in the Output field}

How test will be performed: 

\item{...\\}
    
\end{enumerate}

\subsubsection{Module 2}

...

\subsection{Tests for Nonfunctional Requirements}

\wss{If there is a module that needs to be independently assessed for
  performance, those test cases can go here.  In some projects, planning for
  nonfunctional tests of units will not be that relevant.}

\wss{These tests may involve collecting performance data from previously
  mentioned functional tests.}

\subsubsection{Module ?}
		
\begin{enumerate}

\item{test-id1\\}

Type: \wss{Functional, Dynamic, Manual, Automatic, Static etc. Most will
  be automatic}
					
Initial State: 
					
Input/Condition: 
					
Output/Result: 
					
How test will be performed: 
					
\item{test-id2\\}

Type: Functional, Dynamic, Manual, Static etc.
					
Initial State: 
					
Input: 
					
Output: 
					
How test will be performed: 

\end{enumerate}

\subsubsection{Module ?}

...

\subsection{Traceability Between Test Cases and Modules}

\wss{Provide evidence that all modules have been considered.}
				
\bibliographystyle{plainnat}

\bibliography{../../refs/References}

\newpage

\section{Appendix}

This is where you can place additional information.

\subsection{Symbolic Parameters}

The definition of the test cases will call for SYMBOLIC\_CONSTANTS.
Their values are defined in this section for easy maintenance.

\subsection{Usability Survey Questions?}

\wss{This is a section that would be appropriate for some projects.}

\newpage{}

\section{Appendix --- Generative AI Use Disclosure}

\wss{ChatGPT (version 5) was used to brainstorm ideas for the template for this disclose section.}

\subsection{AI Tools Used}

\wss{Give the name of the AI tool(s) used and the version.}

\subsection{How and Where Used}

\wss{Explain what the tool(s) were used for.  Examples include: brainstorming,
explaining a technical concept, reviewing the document, generating or modifying
text, generating or modifying diagrams, identifying errors or omissions or
inconsistencies.}

\wss{For each of the uses, identify which parts of the document were affected.}

\wss{You do not need to report every interaction, but you should disclose the 
uses that materially contributed to your work}

\subsection{Verification of AI-Generated Content}

\wss{\begin{itemize}
    \item Describe how the team checked AI-generated or AI-assisted material
    before including it in the document.
    \item In particular, verify factual claims, technical assertions,
    requirements, calculations, references, and other information for which an
    AI system may produce incorrect or fabricated results.
    \item \textbf{The team is responsible for the accuracy, completeness,
consistency, and appropriateness of the submitted document, regardless of
whether material was generated by a human or an AI system.}
\end{itemize}}

\newpage{}

\section{Appendix --- Reflection}

\wss{This section is not required for CAS 741}

The information in this section will be used to evaluate the team members on the
graduate attribute of Lifelong Learning.

\input{../Reflection.text}

\begin{enumerate}
  \item What went well while writing this deliverable? 
  \item What pain points did you experience during this deliverable, and how
    did you resolve them?
  \item What knowledge and skills will the team collectively need to acquire to
  successfully complete the verification and validation of your project?
  Examples of possible knowledge and skills include dynamic testing knowledge,
  static testing knowledge, specific tool usage, Valgrind etc.  You should look to
  identify at least one item for each team member.
  \item For each of the knowledge areas and skills identified in the previous
  question, what are at least two approaches to acquiring the knowledge or
  mastering the skill?  Of the identified approaches, which will each team
  member pursue, and why did they make this choice?
\end{enumerate}

\end{document}